% ============================================================================
% CAPITOLO TEMPORANEO — DIARIO DEGLI SVILUPPI
% ----------------------------------------------------------------------------
% Controparte in prosa di ciò che viene realizzato sotto pii_detection/: cosa è
% stato costruito, quali decisioni sono state prese e perché, quali limiti
% restano. È TEMPORANEO: ogni voce è materiale grezzo destinato ai capitoli
% canonici (§5 architettura, §6 implementazione, §7 test) e va RIMOSSA da qui
% una volta travasata nel paragrafo di destinazione, indicato dal rimando.
% Quando il diario è vuoto, eliminare questo file e il relativo \input in
% DOC_Algisi.tex.
% Prosa in italiano; ogni token di codice in inglese dentro \texttt{...}.
% ============================================================================

\chapter{Diario degli sviluppi (\emph{work in progress})}
\label{cap:diario-sviluppi}

\textit{Nota.} Capitolo temporaneo di servizio. Raccoglie, area per area, ciò che è stato realizzato nel sistema con le relative decisioni progettuali, in una forma già utilizzabile come prosa di tesi: ciascuna voce indica il paragrafo canonico in cui confluirà e va rimossa da qui una volta assorbita. Va eliminato a lavoro concluso.

\section{Data di riferimento del documento e provenienza della stima}
\label{sec:sv-reference-date}

\textit{Destinazione: §\ref{sub:b7-retention} (architettura) e il paragrafo sull'estrazione del Capitolo~\ref{cap:implementazione}.} \medskip

La verifica della conservazione confronta l'età di un documento con la durata dichiarata nel registro dei trattamenti, e presuppone quindi di sapere \emph{quando} quel documento risale. Il file system, da solo, è un testimone inaffidabile: la copia di una condivisione di rete, una migrazione o un ripristino da backup riscrivono la data di ultima modifica di ogni file, e un documento del 2019 appare improvvisamente come scritto ieri. Assumere il \emph{mtime} come data del documento significa quindi accettare un falso negativo sistematico proprio nello scenario --- l'archivio storico mai più toccato --- che il controllo dovrebbe intercettare. \medskip

I documenti nativi digitali, però, portano spesso dentro di sé una risposta migliore: il dizionario informativo di un PDF conserva \texttt{modDate} e \texttt{creationDate}, i formati Office espongono le \emph{core properties} con data di creazione e di ultima modifica. Il modulo \texttt{extraction/dates.py} interroga per primo questo livello e ricade sul \emph{mtime} solo quando non trova nulla di utilizzabile, secondo un ordine di preferenza dichiarato: data di modifica interna, data di creazione interna, data del file system. \medskip

La scelta progettuale che qualifica il meccanismo non è però la catena in sé, bensì il fatto che il sistema \textbf{dichiara che cosa ha usato}: il valore restituito (\texttt{ReferenceDate}) porta con sé la provenienza (\texttt{DateSource}, con i due casi \texttt{content\_metadata} e \texttt{file\_mtime}) e il campo esatto da cui la data è stata letta. Il verdetto di conformità può così qualificare quanto vale il segnale che mostra al DPO, invece di presentare una stima come se fosse un accertamento; una segnalazione fondata sul solo \emph{mtime} resta un indizio da verificare, non una violazione conclamata. \medskip

I metadati reali sono peraltro spesso sporchi --- campi inizializzati all'epoca Unix, orologi impostati male, date collocate nel futuro --- e credervi ciecamente reintrodurrebbe l'errore da un'altra porta. Una sola guardia di plausibilità, concentrata in un unico punto del modulo perché non possa divergere da un formato all'altro, scarta i valori anteriori al 1990 o successivi al momento presente e passa al candidato seguente: una data di modifica corrotta accanto a una data di creazione sensata produce quest'ultima, non un ritorno immediato al file system. Allo stesso modo, l'impossibilità di leggere i metadati di un file --- un PDF troncato, proprietà illeggibili --- non è trattata come un errore: datare un documento non deve poter interrompere la sua analisi, e il ritorno al \emph{mtime} è già la risposta corretta. \medskip

Resta un'approssimazione consapevole, ed è opportuno dirlo: i metadati interni possono mancare del tutto (testo semplice, Word legacy, ODF), possono essere stati riscritti da una conversione di formato e possono a loro volta mentire. Il sistema non tenta, allo stato attuale, di ricavare la data dal corpo del documento --- una lettera che si apre con «Milano, 12 marzo 2019» resta illeggibile sotto questo profilo --- perché l'euristica sarebbe rumorosa e andrebbe validata con lo stesso rigore del rilevamento delle PII.

\section{Due orologi distinti nel registro dei documenti}
\label{sec:sv-due-orologi}

\textit{Destinazione: paragrafo sul modello dati del registro, Capitolo~\ref{cap:implementazione}.} \medskip

Il registro delle PII conserva ora, per ciascun documento, \textbf{due} informazioni temporali che è essenziale non confondere, perché rispondono a domande diverse e si comportano in modo diverso. La prima è la \emph{data di riferimento} (\texttt{reference\_date}, con la sua provenienza in \texttt{reference\_date\_source}): è una grandezza \textbf{semantica}, dice quanto è vecchio il \emph{contenuto}, e può provenire dall'interno del file (§\ref{sec:sv-reference-date}); è ciò su cui lavora la verifica della conservazione. La seconda è l'\emph{impronta} del file (\texttt{source\_mtime} e \texttt{source\_size}): è una grandezza \textbf{tecnica}, descrive lo stato osservabile del file su disco al momento dell'ultima analisi, e serve a rispondere a una domanda del tutto diversa --- questo documento è cambiato da quando l'abbiamo guardato, e vale quindi la pena rianalizzarlo? \medskip

La distinzione non è pedanteria terminologica: usare la data semantica anche per la seconda domanda produrrebbe un errore silenzioso. La data di modifica interna di un PDF resta spesso invariata quando il file viene riscritto, e un documento effettivamente modificato apparirebbe intatto, venendo saltato indefinitamente da ogni analisi successiva. Simmetricamente, usare la data del file system come età del contenuto è l'approssimazione che il paragrafo precedente ha messo in discussione. Due domande, due misure, due campi. \medskip

L'impronta è deliberatamente il segnale più economico che funziona --- data di modifica e dimensione, ottenute con una sola interrogazione al file system --- e non un'impronta crittografica del contenuto: calcolare un digest di ogni file di una condivisione aziendale costerebbe esattamente ciò che l'analisi incrementale intende risparmiare. A queste si affianca \texttt{last\_scanned\_at}, il momento dell'ultima analisi effettiva del documento, che una scansione in cui il documento venga saltato non aggiorna. \medskip

Ne discende una regola di aggiornamento precisa: date e impronta si rinfrescano \emph{solo} quando la scansione ha effettivamente osservato il file. La riconciliazione dei documenti spariti dal disco, che registra una scansione vuota per marcarne le PII come rimosse, non ha alcun file da guardare e non deve perciò azzerare ciò che l'ultima osservazione reale aveva registrato.

\section{Il silenzio non è conformità: conservazione non verificabile}
\label{sec:sv-retention-non-verificabile}

\textit{Destinazione: §\ref{sub:b7-retention} (architettura) e il paragrafo sul verdetto di conformità, Capitolo~\ref{cap:implementazione}.} \medskip

Il controllo di conservazione confronta l'età del documento con la durata dichiarata nel registro, ma quella durata non è sempre un numero. I registri reali esprimono spesso un \emph{criterio} anziché un termine --- «per la durata del rapporto di lavoro», «fino alla revoca del consenso» --- e l'ingestione, non trovando alcuna quantità da estrarre, lascia la durata indefinita. Nella prima versione del controllo quelle categorie venivano semplicemente saltate, con una conseguenza che è opportuno enunciare per quello che è: il documento risultava conforme sull'asse della conservazione senza che alcuna verifica fosse mai avvenuta. Il caso «nessuno ha potuto controllare» era indistinguibile dal caso «controllato e in regola». \medskip

Il verdetto distingue ora esplicitamente i due esiti. Ogni macro-categoria dichiarata i cui dati siano ancora presenti nel documento produce, in alternativa, una \textbf{violazione} --- quando il documento è più vecchio del termine dichiarato --- oppure una segnalazione di \textbf{conservazione non verificabile}, quando il registro esprime un criterio. La seconda non rende il documento non conforme, esattamente come la segnalazione delle categorie dichiarate ma non riconducibili ad alcun tipo rilevabile: è un'informazione destinata al DPO, che indica dove il registro andrebbe reso più preciso. Vale la pena osservare che questa segnalazione non dipende dalla data del documento: la lacuna è nel registro, non nel documento, e va comunicata anche quando la data non è ricostruibile. \medskip

Alla stessa logica appartiene un'altra asimmetria corretta in questa fase. Il confronto fra dichiarato e rilevato considera, di norma, le sole categorie \emph{confermate} dal DPO, ignorando le associazioni ancora in stato di proposta; il controllo di conservazione, invece, le contava tutte. Lo stesso documento poteva così ricevere due giudizi incoerenti sulla medesima categoria --- «non dichiarata», e quindi orfana, per un ramo del verdetto; «dichiarata e conservata oltre il termine» per l'altro. La regola su quali categorie contino è ora scritta in un solo punto del codice e utilizzata da entrambi i rami, così che i due non possano più divergere. \medskip

Infine, sul versante dell'ingestione, il riconoscimento delle durate è stato esteso alle unità francesi (\texttt{an}, \texttt{ans}, \texttt{mois}) accanto a quelle inglesi e italiane già supportate. Non è un dettaglio linguistico: il modello di registro adottato come riferimento è quello dell'autorità francese, e una dicitura come «5 ans» lasciata non interpretata avrebbe disattivato in silenzio il controllo di conservazione per quella categoria --- ricadendo esattamente nel difetto descritto sopra, ma per una causa evitabile.


\section{Dalla scheda del documento alla vista d'insieme}
\label{sec:sv-vista-conservazione}

\textit{Destinazione: paragrafo sulla dashboard del DPO, Capitolo~\ref{cap:implementazione}.} \medskip

Il verdetto di conformità nasce come giudizio su un singolo documento, ed è la forma corretta per la domanda «questo file è in regola?». Non è però la domanda che un DPO rivolge a una condivisione aziendale: la sua è «mostrami tutto ciò che è oltre il termine». Con trecento documenti, ottenere quella risposta aprendo una scheda alla volta non è una procedura praticabile, ed è il modo più sicuro perché una violazione reale resti invisibile in mezzo al rumore. \medskip

È stata perciò aggiunta una vista d'insieme, disponibile sia come pagina web sia come comando da terminale, che applica il medesimo verdetto all'intero registro e ne conserva soltanto le righe che hanno qualcosa da dire, ordinate per \textbf{gravità}, cioè per il numero di mesi di sforamento. La misura della gravità non è un dettaglio estetico: un mese oltre il termine è un promemoria, dieci anni oltre sono un incidente, e un elenco che non li distingue costringe a rileggerlo tutto. \medskip

Tre vincoli hanno guidato la realizzazione. Primo, \textbf{nessuna logica nuova}: la vista richiama la stessa funzione che produce il verdetto del singolo documento, così le due schermate non possono divergere. Secondo, \textbf{nessuna scrittura}: consultare una pagina non deve modificare lo stato del sistema, mentre il controllo per-documento, quando invocato per l'assegnazione, persiste la copertura di ciascuna PII. Terzo, il registro dei trattamenti viene letto \textbf{una sola volta} e tenuto in memoria per la durata della vista, anziché una volta per documento. \medskip

I documenti privi di associazione a un trattamento sono deliberatamente esclusi: senza un trattamento non esiste una durata dichiarata con cui confrontarsi, e includerli significherebbe sommergere le violazioni effettive sotto ogni file non ancora classificato della condivisione. È una scelta che vale la pena rendere esplicita nella pagina stessa, perché un elenco vuoto non significhi «tutto a posto» quando in realtà significa «non c'è ancora nulla da confrontare».

\section{Rianalizzare solo ciò che è cambiato}
\label{sec:sv-scansione-incrementale}

\textit{Destinazione: §\ref{sub:rianalisi-incrementale} (requisiti) e il paragrafo sulla scansione batch, Capitolo~\ref{cap:implementazione}.} \medskip

La scansione di una cartella rileggeva per intero, a ogni esecuzione, ogni documento incontrato: estrazione del testo e rilevamento su file identici a quelli della volta precedente. Su un albero aziendale è il costo dominante dell'intera procedura, ed è un costo interamente sprecato; soprattutto, rende impraticabile l'esecuzione periodica che il modello d'impiego previsto --- «si indica una cartella e il sistema se ne occupa» --- presuppone. Il registro sapeva già \emph{che cosa} fosse cambiato nel contenuto di un documento, grazie al confronto fra scansioni successive, ma non sapeva \emph{se valesse la pena guardarlo}. \medskip

Il rimedio adottato è il più semplice che funzioni: si confrontano la data di modifica e la dimensione del file con quelle registrate all'ultima analisi e, se coincidono, il documento non viene riletto. Si è scelto deliberatamente di non calcolare un'impronta crittografica del contenuto, che sarebbe più precisa ma richiederebbe di leggere ogni byte di ogni file --- esattamente il costo che la misura intende evitare. Il confronto è per \emph{differenza} e non per posteriorità: un file ripristinato a una versione precedente è cambiato quanto uno modificato, e un criterio del tipo «più recente di» lo escluderebbe per sempre dall'analisi. Una tolleranza di un secondo assorbe le diverse granularità temporali dei file system, che altrimenti farebbero apparire modificato ogni documento. \medskip

Il punto delicato dell'intervento non è però il criterio, bensì un'invariante preesistente. La scansione riconcilia il registro con il disco marcando come rimosse le PII dei documenti spariti, e per farlo si basa sull'insieme dei documenti «visti». Se un documento saltato perché invariato non entrasse in quell'insieme, la seconda scansione dichiarerebbe rimosse le PII di ogni file mai toccato --- un danno molto peggiore del costo che si voleva risparmiare. Il significato dell'insieme è stato quindi reso esplicito: vi appartiene ogni file \textbf{trovato sul disco}, indipendentemente da ciò che è accaduto dopo. La stessa precisazione ha corretto un difetto preesistente e indipendente: un file che fallisce l'estrazione non entrava in quell'insieme, e le sue PII venivano dichiarate rimosse pur essendo il file ancora al suo posto. Coerentemente, un file mai analizzato con successo non viene mai considerato invariato, ma ritentato a ogni esecuzione. \medskip

Un'ultima insidia riguarda il motore di rilevamento anziché i documenti. Abilitare un modello di riconoscimento più capace, o aggiungere una regola alla configurazione dei pattern, cambia ciò che verrebbe trovato in un file i cui byte non si sono mossi: il registro continuerebbe a mostrare i risultati del motore precedente presentandoli come aggiornati. Un registro fermo è un problema; un registro che mente sul proprio aggiornamento è peggio. Per questo, accanto all'impronta del file, viene registrata un'impronta della \textbf{configurazione di rilevamento} --- identificatori dei rilevatori attivi e contenuto dei file di configurazione --- e una sua variazione invalida da sola il salto, senza che l'operatore debba ricordarsi di forzare una rianalisi completa. \medskip

Due comportamenti restano volutamente fuori dall'ottimizzazione. L'anteprima della scansione dichiara in anticipo quanti file verrebbero saltati, perché la pagina che promette di dire «che cosa verrà analizzato» deve continuare a dire il vero anche adesso che la risposta è più articolata. E i file caricati dal browser vengono sempre rianalizzati per intero: materializzati sul server al momento del caricamento, portano come data di modifica quella del caricamento stesso, sicché nessuna impronta calcolata su di essi significherebbe qualcosa.

\section{Rilevamento generativo come secondo parere}
\label{sec:sv-ai-detector}

\textit{Destinazione: §\ref{sec:scelta-modello-llm} e §\ref{sub:slm-task} (studi, per l'inquadramento teorico e la scelta del modello), §\ref{sec:impl-detection} (implementazione, per la realizzazione del rilevatore) e il Capitolo~\ref{cap:studi} per l'esito comparativo del benchmark.} \medskip

\textbf{Perché questa funzionalità esiste.} L'intero apparato di rilevamento costruito finora --- espressioni regolari con checksum, riconoscimento di entità zero-shot, fusione dei due --- appartiene al paradigma degli algoritmi \emph{tradizionali}: veloci, deterministici, economici, ma ciechi a tutto ciò che non è stato previsto in una regola o in un'etichetta. La domanda che la tesi si pone --- ed è una domanda di merito, non di ingegneria --- è se e quando un modello generativo, che legge il documento come farebbe un lettore umano, trovi dato personale che quegli algoritmi non trovano, e se il guadagno giustifichi il suo costo, che è di ordini di grandezza superiore in tempo ed energia. Rispondere richiede di poter \emph{misurare} l'AI sullo stesso terreno degli altri rilevatori; e per misurarla sullo stesso terreno, l'AI deve essere costruita come un rilevatore fra gli altri, non come un meccanismo a parte. \medskip

\textbf{Un rilevatore, non un'eccezione.} Il modello generativo è stato perciò incapsulato dietro lo stesso contratto (\texttt{PIIDetector}) che espongono il rilevatore a regole e quello neurale: riceve il testo di un documento e restituisce candidati con posizione, tipo e provenienza. Questa uniformità non è un vezzo architetturale: è la condizione che rende l'AI confrontabile: la stessa fusione che oggi concilia regole e NER può accogliere l'AI come terza voce, e lo stesso banco di prova che assegna precisione e richiamo a un rilevatore può assegnarli al modello generativo senza codice speciale. L'architettura riservava da tempo a questa tecnica un posto già nominato --- una categoria \texttt{AI} fra i tipi di rilevatore, un livello di conferma \texttt{AI\_DISCOVERED}, un parametro nella fusione destinato ai candidati dell'AI --- ma quel posto era rimasto vuoto: questa fase lo riempie. \medskip

\textbf{Il problema del posizionamento e la difesa dalle allucinazioni.} Un modello linguistico sa dire \emph{quali} valori sono dato personale, ma non \emph{dove} si trovino nel testo: gli offset che eventualmente produce sono inaffidabili, perché il modello non conta caratteri. La soluzione adottata rovescia il problema: al modello si chiedono i soli \emph{valori} (la stringa esatta, copiata dal testo), e il posizionamento lo esegue il sistema, cercando ogni valore restituito dentro il testo reale e registrandone tutte le occorrenze. Questo scarto fra ciò che il modello afferma e ciò che il testo conferma diventa, non secondariamente, la prima difesa contro le \emph{allucinazioni}: un valore che il modello inventa e che nel documento non compare non viene semplicemente trovato, e cade. A questa si affianca la seconda difesa, mutuata dal mappatore di categorie del registro: il modello riceve nel prompt il \textbf{catalogo chiuso} dei tipi ammessi e può assegnare solo quelli; ogni tipo fuori catalogo è scartato. Il modello, in altre parole, non è mai creduto sulla parola: è creduto solo dove il testo e il catalogo lo confermano. \medskip

\textbf{Un potenziamento, mai un punto di rottura.} Il modello gira su un runtime locale che può essere spento, lento o assente. Il rilevatore ne tiene conto per costruzione: analizza il documento a finestre sovrapposte --- così un dato a cavallo di due finestre resta comunque leggibile per intero da una delle due --- e il fallimento di una finestra, o l'irraggiungibilità del runtime, non interrompe l'analisi ma la degrada, restituendo i risultati delle finestre riuscite o, al limite, nessun risultato. Un rilevatore che può fare crollare l'intera scansione non è un potenziamento: è una fragilità. La certezza attribuita a un rilevamento dell'AI è deliberatamente moderata --- inferiore a quella di un'espressione regolare validata da un checksum, superiore a quella del nulla --- perché un'affermazione singola e non verificata di un modello vale meno di una regola che ha superato una prova aritmetica; sarà la fusione, quando l'AI e un altro rilevatore concordano sullo stesso dato, ad alzarla. \medskip

\textbf{La fusione: conferma e scoperta.} È nella fusione che il secondo parere dell'AI acquista significato, e il modo in cui vi entra è ciò che rende \emph{misurabile} l'accordo fra generativo e tradizionale --- l'informazione che alla tesi interessa più di ogni altra. La stessa fusione che concilia regole e riconoscimento di entità accoglie ora i candidati dell'AI per ultimi, appaiando ciascuno al rilevamento con cui si sovrappone di più, con la medesima soglia geometrica usata fra gli altri due rilevatori. Da qui l'AI gioca due ruoli. Quando ricade su un dato già trovato e ne concorda il tipo è \textbf{conferma}: la sua provenienza si aggiunge alle altre e la certezza del rilevamento sale; un dato prima sostenuto da una sola voce diventa a doppia conferma, uno già a doppia conferma acquista una terza voce. Quando invece cade dove nessun altro rilevatore aveva visto nulla è \textbf{scoperta}: un rilevamento a sé, marcato come trovato dalla sola AI, che è esattamente il caso su cui la tesi misura se il modello generativo veda ciò che gli algoritmi non vedono. \medskip

\textbf{Recall-first anche nel disaccordo.} Restava da decidere il caso scomodo: l'AI si sovrappone a un rilevamento ma ne afferma un tipo diverso. La scelta adottata non arbitra e non scarta --- coerente con il principio di privilegiare il richiamo e di rimandare ogni arbitrato al livello di persistenza --- ma pesa le due voci in disaccordo per quanto valgono. Un parere singolo dell'AI che contraddice un dato sostenuto da un solo altro rilevatore mette i due sullo stesso piano, e li segnala entrambi come in conflitto; ma un parere singolo dell'AI non basta a incrinare un accordo già raggiunto fra regole e riconoscimento di entità, che resta quindi tale, con la voce discorde tenuta accanto come rilevamento a sé. La certezza, in altre parole, non si perde per una contestazione più debole di ciò che contesta. Una piccola guardia completa il quadro: uno stesso rilevatore dell'AI che, per effetto delle finestre sovrapposte, riproponesse due volte lo stesso dato non ne gonfia la certezza, perché una provenienza già presente non viene conteggiata una seconda volta. \medskip

\textbf{Quando accendere l'AI: una politica del chiamante.} Poiché una passata generativa costa da secondi a minuti per documento, non la si esegue su ogni file a ogni scansione. La decisione su \emph{quali} documenti sottoporre all'AI è però tenuta fuori dal rilevatore --- che si limita ad analizzare il testo che riceve --- e affidata a una \textbf{politica} del chiamante: un campionamento di un documento ogni $N$, deterministico e privo di stato, che non dipende dalla storia delle scansioni e sopravvive a un riavvio. Tenere separata la politica dal rilevatore è ciò che consente di innescare l'AI in modi diversi --- a campione durante una scansione, oppure su richiesta su un singolo documento --- senza duplicare la logica di rilevamento. Nel dispiegamento containerizzato il campionamento è attivo per default a un documento ogni dieci --- un compromesso fra copertura e costo che l'operatore alza, abbassa o azzera da una sola variabile d'ambiente --- mentre il modello del rilevatore è configurato a parte da quello del mappatore di categorie, così che i due compiti generativi possano usare pesi diversi. \medskip

\textbf{L'innesto: uno slot opzionale, non un registro.} Collegare il nuovo rilevatore alle vie con cui il sistema è realmente usato --- l'analisi di un singolo documento, l'ingestione nel registro, la scansione ricorsiva di una cartella --- poteva richiedere un'infrastruttura di orchestrazione dei rilevatori; si è preferita la strada più semplice e meno ripetitiva, coerente con l'uniformità di contratto già scelta. L'AI è un \emph{parametro opzionale} che attraversa quelle tre funzioni: assente, la pipeline si comporta esattamente come prima e nessuno dei percorsi esistenti cambia; presente, i suoi candidati entrano nella fusione. La scelta si legge nel vocabolario dei pattern di progettazione: i rilevatori sono \emph{Strategy} intercambiabili dietro un contratto comune (il \texttt{Protocol} \texttt{PIIDetector}) e l'AI vi entra per \emph{dependency injection}, cosicché aggiungerla è un'\emph{estensione} --- una nuova strategia e un parametro --- e non una \emph{modifica} del codice esistente (principio Open/Closed). È proprio questo a rendere superfluo un pattern di orchestrazione più pesante --- un registro dei rilevatori, un mediatore: la stessa estensibilità che quel pattern offrirebbe è già data dalla combinazione Strategy più iniezione, e introdurlo sarebbe complessità non ripagata. Il campionamento prende corpo nella scansione di cartella, dove a ciascun documento è associato il suo indice nell'\emph{intera} enumerazione dell'albero --- non nella sola parte effettivamente riletta --- così che il sottoinsieme campionato resti lo stesso indipendentemente da quali file la scansione incrementale abbia saltato; la decisione «quali documenti» è a sua volta un oggetto a sé (una \emph{policy}), separato dal rilevatore. La stessa impronta del motore di rilevamento, che decide cosa rileggere, ignora di proposito l'AI: essa è un parere occasionale e campionato, non parte del motore deterministico, e non deve perciò forzare la rilettura di ciò che non è cambiato. Da riga di comando la funzionalità è già interamente percorribile --- un'opzione dedicata attiva l'AI su un documento o su un'intera cartella, e il campionamento periodico si accende da una variabile d'ambiente --- il che chiude la prima dimostrazione completa, dall'estrazione al verdetto, con il secondo parere generativo incluso. \medskip

\textbf{Rendere visibile la conferma nel registro.} Un dettaglio della persistenza avrebbe altrimenti reso muto proprio l'esito più interessante. Quando una ri-scansione ritrova una PII nella stessa posizione, il registro la tratta come \emph{confermata} e finora si limitava ad aggiornarne la data dell'ultima osservazione, non la certezza. Ma è esattamente ciò che accade quando l'AI, campionata su quel documento a una scansione successiva, concorda su un dato che prima era sostenuto da una sola voce: la posizione non cambia --- dunque «confermata», non «spostata» --- eppure il livello di conferma, la fiducia e l'elenco delle provenienze sì. Il ramo di conferma rinfresca ora quei tre campi come già faceva quello di spostamento, così che l'accordo sopraggiunto diventi leggibile; l'identità del dato (tipo e posizione) e la natura dell'evento registrato restano quelle di una conferma. \medskip

\textbf{L'AI nell'interfaccia: un'analisi che non blocca.} Perché la funzionalità sia realmente usabile dal DPO, e non solo da riga di comando, l'AI compare nei due punti in cui il sistema è operato. All'avvio di una scansione, un interruttore la estende a tutti i documenti dell'albero; in assenza di quell'interruttore, il campionamento periodico governato dall'ambiente resta comunque attivo, e la pagina di stato dichiara quale dei due regimi è in corso --- «su tutti» oppure «campionati, uno ogni $N$». Il secondo punto è la scheda del singolo documento, con un'azione di rianalisi su richiesta, ed è qui che una scelta di progettazione era obbligata. L'analisi generativa di un documento lungo su CPU dura minuti: eseguirla dentro la richiesta web terrebbe l'utente bloccato su una pagina che gira, e renderebbe il sistema inservibile fino al termine. La rianalisi è perciò \textbf{asincrona}. Non è però stato costruito nulla di nuovo per ottenerla: si è \emph{riusato lo stesso meccanismo di job in background} già esistente per la scansione di cartelle --- un'esecuzione su un thread separato, registrata in memoria, di cui la pagina interroga periodicamente lo stato --- applicandolo a un solo file. L'azione avvia il job e rimanda subito a una pagina di avanzamento; l'utente è libero di navigare altrove e tornare quando vuole, e a lavoro concluso la scheda del documento mostra le nuove istanze confermate o scoperte dall'AI, senza che la resa di quella pagina abbia richiesto alcuna modifica --- livello di conferma e provenienze erano già ciò che essa mostrava. Riutilizzare un meccanismo maturo invece di scrivere un secondo percorso asincrono è, di nuovo, la scelta che costa meno ed è meno esposta a errori. \medskip

\textbf{Dalla PII orfana alla proposta di registro.} La ragione per cui si fa scoprire dato personale dall'AI non è compilare un inventario, ma confrontarlo con ciò che il registro dei trattamenti dichiara: una PII trovata e non dichiarata da alcun trattamento è già, nel verdetto di conformità, una \emph{orfana} --- il rischio nella sua forma più pura. Finora quel rischio era però solo \emph{segnalato}. L'ultimo tassello lo rende \emph{azionabile}: dalla scheda del documento, dove l'orfana è già in evidenza, il DPO può trasformarla in una \textbf{proposta} di categoria da aggiungere a uno dei trattamenti associati. La proposta nasce in uno stato esplicito di «proposta», non di categoria confermata, e vive sotto un raggruppamento dedicato \emph{senza durata di conservazione}: il sistema porta \emph{che cosa} è stato trovato e \emph{dove}, mai una conservazione o una base giuridica che non conosce. La conferma vera resta un atto del DPO nella revisione del registro; il sistema non scrive mai da sé la \emph{ground truth} dichiarativa, la propone soltanto. Una scelta merita di essere resa esplicita: il nome della categoria proposta è l'etichetta che il catalogo dei tipi già associa a quel \texttt{pii\_type}, presa così com'è, e non un testo chiesto al modello generativo. Nel solo punto in cui il flusso si avvicina alla scrittura del registro si è preferito un dato deterministico a un'inferenza che potrebbe allucinare: il catalogo quel nome ce l'ha già, non c'è ragione di farlo inventare. Sul piano del verdetto la proposta si chiude da sé: essendo una categoria dichiarata, alla successiva verifica la stessa PII risulta coperta anziché orfana, e --- non avendo durata --- compare fra le conservazioni «non verificabili», che è esattamente ciò che è, finché il DPO non la completa. \medskip

\textbf{Il banco di prova: la domanda di tesi in una tabella.} Tutto ciò che precede predispone il rilevatore; questo passo lo \emph{misura}, ed è il motivo per cui la funzionalità esiste. Sullo stesso corpus annotato, per ciascun modello candidato, si misurano tre grandezze affiancate: la qualità (precisione, richiamo, F1), la latenza (secondi per documento) e l'energia (wattora per documento). Tre tipi di riga rendono visibile il contributo \emph{netto} dell'AI: una riga di riferimento --- il sistema attuale, regole più NER fuse, senza AI; una riga per l'AI \emph{da sola}, per ciascun modello; e una riga di \emph{unione} --- il sistema con quel modello. Lo scarto fra unione e riferimento è ciò che l'AI aggiunge; la riga dell'AI da sola dice quanto il modello trova per proprio conto. \medskip

\textbf{Perché più fasce di taglia, e non solo la più piccola.} L'elenco predefinito dei modelli attraversa deliberatamente tre fasce di grandezza --- circa quattro miliardi di parametri, sette-otto, dodici-quattordici. Limitare la prova alla fascia più piccola misurerebbe il \emph{compromesso di dispiegamento} (l'unico peso che gira comodo ovunque), non il \emph{potenziale} del modello generativo, che è la domanda di merito: quanto recupera un modello più grande sulle PII contestuali che gli algoritmi non trovano, e vale i secondi e i wattora in più che costa su CPU? La tabella è ordinata per taglia crescente proprio perché quel compromesso si legga scorrendola. La misura dell'energia va presa per ciò che è --- una \emph{stima comparativa}, non un contatore: dove il conteggio hardware non è disponibile, la libreria adottata ripiega su una stima dal profilo di consumo del processore, e l'inferenza gira comunque nel container del modello, accanto e non dentro il processo dell'applicazione. I numeri hanno senso gli uni rispetto agli altri, non in assoluto, e vanno prodotti eseguendo la prova in container, dove i modelli girano; per questo l'apparato è costruito e verificato con rilevatori fittizi, ma la tabella di esito su modelli reali è un esperimento da condurre, non un dato già in mano. \medskip

\textbf{Minimizzazione.} Il testo dei documenti, valori delle PII compresi, esce dal processo dell'applicazione soltanto verso il modello \emph{locale}, mai verso un servizio esterno; è una precisazione dovuta in uno strumento che esiste per ridurre l'esposizione del dato. Al modello si può inoltre chiedere una breve \emph{motivazione} di ciascun rilevamento, utile in fase di ispezione, ma poiché tale motivazione può citare il valore trovato essa resta in memoria per la durata dell'analisi e non viene mai scritta nel registro, coerentemente con il vincolo di non persistere i valori. \medskip

\textbf{Che cosa manca ancora.} Il rilevatore, la sua fusione, l'innesto nelle vie da riga di comando, la presenza nell'interfaccia web, la proposta di arricchimento del registro e l'apparato di misura sono tutti realizzati, con il comportamento fissato da prove che non richiedono un modello in esecuzione. Resta una cosa sola, e non è codice: \emph{eseguire} il banco di prova sui modelli reali in container e riportarne l'esito --- l'esperimento che dà alla funzionalità la sua risposta di merito. \medskip
